\SourceLeaf{276}%
\noindent Die Anwendung dieser Reihe auf die oben betrachteten einzelnen Fälle gibt
\[
\begin{aligned}
\text{für}\quad p&=-6,\quad\log m=+0{,}000001050448\\
\text{für}\quad p&=+6,\quad\log m=-0{,}000001096531
\end{aligned}
\]
Die endliche Formel 4, welche man auch so schreiben kann
\[
\begin{aligned}
m&=\frac{\alpha A\cos(Q+q)\BareRoot{(1-ee\sin(P+p)^2)}}{a\cos(P+p)}\\
&=\frac{\cos\eta\cos(Q+q)\BareRoot{(1-ee\sin(P+p)^2)}}{\cos\theta\cos(P+p)}
\end{aligned}
\]
gibt, mit zehnziffrigen Logarithmen berechnet, bis auf die zehnte Ziffer genau dasselbe.
\Article{8}
\indent Für die umgekehrte Aufgabe, wo \(q\) gegeben und \(p\) gesucht wird, ist die Entwickelung in eine Reihe noch wesentlicher, da die endliche Formel 3 in diesem Falle nur auf indirectem Wege zum Ziele führen könnte. Der \textsc{Taylor}sche Lehrsatz gibt
\[
p=\frac{\mathrm{d}w}{\mathrm{d}U}\mathord{.}q-\frac{\mathrm{dd}w}{2\mathrm{d}U^2}\mathord{.}qq+\frac{\mathrm{d}^3w}{6\mathrm{d}U^3}\mathord{.}q^3-\text{u.\,s.\,f.}
\]
wo für die Differentialquotienten diejenigen bestimmten Werthe zu setzen sind, welche zu \(q=0\) oder \(U=90^0-Q\), \(w=90^0-P\) gehören. Für die unbestimmten Werthe der drei ersten Differentialquotienten ergeben sich folgende Ausdrücke
\[
\frac{\mathrm{d}w}{\mathrm{d}U}=\frac{(1-ee\cos w^2)\sin w}{\alpha(1-ee)\sin U}
\]
\[
\begin{aligned}
\frac{\mathrm{dd}w}{\mathrm{d}U^2}={}&\frac{(1-ee\cos w^2)\sin w}{\alpha\alpha(1-ee)^2\sin U^2}\\
&\quad(\alpha(1-ee)\cos U-\cos w+ee\cos w(\cos w^2-2\sin w^2))
\end{aligned}
\]
\[
\begin{aligned}
\frac{\mathrm{d}^3w}{\mathrm{d}U^3}={}&\frac{(1-ee\cos w^2)\sin w}{\alpha^3(1-ee)^3\sin U^3}\{\alpha\alpha(1-ee)^2(\cos U^2+2\sin U^2)\\
&-3\alpha(1-ee)\cos U\cos w(1-ee(\cos w^2-2\sin w^2))\\
&+\cos w^2-\sin w^2-ee(2\cos w^4-12\cos w^2\sin w^2+2\sin w^4)\\
&+e^4(\cos w^6-11\cos w^4\sin w^2+6\cos w^2\sin w^4)\}
\end{aligned}
\]
\indent Die beiden folgenden gleichfalls vollständig entwickelten Coëfficienten setze ich um den Raum zu schonen, nicht hieher, da sie doch nur Zwischengrössen sind, um zu den Endresultaten zu gelangen. Diese finden sich nach der Sub%
